\section{项目实践收获}
\subsection{项目实践收获}
通过本次课程大作业，我们系统体验了知识图谱从语料处理到应用展示的完整工程流程。相比只停留在单一算法实现，本项目让我们更加认识到知识图谱系统本质上是一项系统工程，具体收获如下：

\begin{itemize}[leftmargin=2em, itemsep=0.5ex]
    \item \textbf{数据治理层面：} 建立并强化了“数据治理先于知识建模”的工程意识。实践证明，PDF 文本抽取、乱码处理与标准化清洗才是决定后续抽取上限的基础。
    \item \textbf{知识建模层面：} 体会到领域本体设计是一个在“语义表达能力”与“工程可实现性”之间动态权衡的过程，训练了我们在复杂任务中进行抽象建模与粒度控制的能力。
    \item \textbf{团队协作层面：} 深入理解了工程任务中模块化拆分、接口对齐与并行开发的重要性，双人协作模式有效提高了联调效率与问题定位速度。
    \item \textbf{系统实现层面：} 分层流水线的设计显著提升了系统的可维护性与可复现性，使我们学会了如何构建符合真实项目思路的可迭代工程架构。
    \item \textbf{实验评估层面：} 认识到从“做出系统”走向“证明有效”的重要性。只有通过构建金标准并进行定量分析，才能真正明确一个图谱系统的性能边界。
\end{itemize}

\subsection{项目特色与创新点}
本项目的特色与创新主要体现在以下几个方面：
\begin{enumerate}[leftmargin=2.2em,label=\arabic*.]
    \item \textbf{完整的工程闭环。} 项目并非只关注某一个局部算法，而是实现了从文献获取、文本预处理、术语抽取、实体关系抽取、图谱构建、多格式导出到前端可视化和算法评估的完整闭环。这种从“数据—算法—系统—实验”全链条展开的组织方式，使项目更接近实际工程原型，而不是停留在单点功能演示层面。
    \item \textbf{较强的规范性。} 项目采用模块化目录结构和标准化中间文件格式，统一使用 \texttt{jsonl}、\texttt{csv}、\texttt{json} 和 \texttt{ttl} 等格式保存中间结果与最终结果。同时引入人工金标准、Precision、Recall 和 F1 等标准评价方式，对知识抽取效果进行量化分析，从而保证整个系统具备较好的可复现性与可解释性。
    \item \textbf{“全自动发现 + 人工介入优化”双模混合知识构建方法。} 本项目提出并实现了这一双轨式知识构建方案。一方面，系统支持基于统计特征（TF-IDF/PMI）与 LLM 语义精筛的全自动知识发现流水线，能够快速处理海量无标注文献；另一方面，系统能够无缝检测并融合人工构建的实体词典与金标准三元组，并利用人工标注集对自动模型进行动态评估与针对性优化。这种方案既保留了大规模抽取的高召回率，又通过人工金标准的引入确保了核心知识的高度精准，充分体现了“机器发现、人工辅助、指标驱动”的现代知识图谱构建思想。
    \item \textbf{面向中文 PDF 的容错处理流程。} 针对中文学术 PDF 常见的乱码、文本层异常和扫描页问题，项目设计了“文本层抽取—细粒度回退—OCR 兜底”的多级容错处理流程；针对关系抽取易产生伪三元组的问题，引入触发词位置、实体黑名单、距离阈值和类型约束四重过滤机制。这些设计直接面向真实文档噪声和抽取误差来源，具有较强工程针对性。
    \item \textbf{轻量化可视化展示方案。} 项目构建了基于 FastAPI 与 ECharts 的轻量化知识图谱可视化系统，无需依赖复杂图数据库即可完成图谱展示、搜索、筛选和详情查看。该实现方式显著降低了部署与演示门槛。
\end{enumerate}
综合来看，本项目的特色与创新并不体现在某一个单独模型上，而体现在面对真实领域文献时所形成的一套完整、规范、可展示、可扩展的知识图谱原型方案。

\subsection{项目不足与改进方向}
尽管本项目已经完成了较完整的知识图谱原型系统，但从更高标准的研究和工程应用要求来看，仍然存在若干不足。首先，在实体识别方面，当前方法仍以词典匹配和规则补充为主，对新术语自动发现能力有限。当文献中出现未被术语表覆盖的新概念、缩写变体或上下文依赖较强的表达时，系统可能出现漏识别或类型判断不稳定的问题。尤其对于嵌套实体、长复合术语和带参数修饰的短语，现有方法仍有进一步优化空间。

其次，在关系抽取方面，虽然系统已经引入了大语言模型辅助处理复杂长难句中的隐式关联，但在极端长文本跨段落的推理上仍有局限。变构飞行器领域文献中经常出现跨段落嵌套、长篇并列列举和分散于不同章节的多实体逻辑共现。目前的并发批处理策略容易产生上下文截断，导致超出单次批处理窗口的因果逻辑链条断裂。因此，未来应考虑引入 RAG（检索增强生成）技术、图神经网络推理算法或更长上下文窗口的专用模型，以实现跨篇章级别的深度关系发现。

在系统应用层面，目前前端已具备基本的图谱展示、搜索和筛选能力，但尚未进一步扩展到路径分析、语义问答、关系推荐和知识演化分析等更高层次功能。若后续将图谱与图数据库、问答系统或推理模块结合，系统可从“展示型原型”逐步演进为“分析型工具”。此外，当前图谱更新流程仍以离线批处理为主，若未来加入增量更新机制与版本管理能力，将更有利于长期维护和持续扩充。

综合来看，本项目的后续改进方向可以概括为两个层面：一是增强抽取层的上下文理解能力，提高对复杂实体和复杂关系的处理效果；二是拓展应用层的交互分析与智能服务能力，使知识图谱从课程级展示系统逐步走向更具实用价值的领域知识服务平台。

\subsection{小组成员分工明细}
本项目由两人小组共同完成，双方在算法设计、系统开发与报告撰写上紧密配合，工作量均等。具体分工明细如表~\ref{tab:division} 所示。

\begin{table}[H]
\centering
\caption{小组成员分工明细}
\label{tab:division}
\renewcommand{\arraystretch}{1.3}
\resizebox{\textwidth}{!}{%
\begin{tabular}{>{\centering\arraybackslash}p{2.5cm}>{\centering\arraybackslash}p{3cm}>{\raggedright\arraybackslash}p{8cm}>{\centering\arraybackslash}p{2cm}}
\toprule
姓名 & 学号 & 负责核心工作内容 & 贡献比例 \\
\midrule
杨博文 & 3023244317 & $\bullet$ 需求分析与技术方案设计 \newline $\bullet$ 领域文献收集与文本清洗 \newline $\bullet$ 术语提取（统计特征 + LLM 精筛） \newline $\bullet$ 实体关系抽取规则设计 \newline $\bullet$ 人工金标准构建与效果评估 \newline $\bullet$ 技术报告联合撰写与修订 & 50\% \\
\midrule
于露雪 & 3023244303 & $\bullet$ 系统架构设计与工程部署 \newline $\bullet$ PDF 文本提取与容错模块 \newline $\bullet$ 大模型并行抽取与词典增量逻辑 \newline $\bullet$ FastAPI 后端开发与图谱多格式导出 \newline $\bullet$ ECharts 前端交互系统开发 \newline $\bullet$ 技术报告主笔与排版 & 50\% \\
\bottomrule
\end{tabular}%
}
\end{table}

\subsection{总结}
本文围绕“变构飞行器领域知识图谱构建与可视化系统设计与实现”这一题目，完成了领域数据整理、文本预处理、术语构建、实体关系抽取、图谱导出、可视化展示和评估分析等关键任务。整体上，本项目已经实现了一个具有较好完整性、规范性和一定创新性的课程原型系统。
